% Part: first-order-logic
% Chapter: axiomatic-deduction
% Section: provability-quantifiers

% verification of properties of provability needed for maximally
% consistent sets in the completeness chapter.

\documentclass[../../../include/open-logic-section]{subfiles}

\begin{document}

\olfileid{fol}{axd}{qpr}

\olsection{\usetoken{S}{derivability} والمكمّمات}

\begin{explain}
  ومن مقدمات برهان الاكتمال أن تثبت بالاشتقاقات البديهية هذه
  النتائج في علاقة~$\Proves$؛ فنبيّنها الآن.
\end{explain}

\begin{thm}
\ollabel{thm:strong-generalization} إذا كان $c$ !!a{constant} لا يرد
في $\Gamma$ ولا في $!A(x)$، وكان $\Gamma \Proves !A(c)$، فإن $\Gamma
\Proves \lforall[x][!A(x)]$.
\end{thm}

\begin{proof}
% تصحيح برهاني (0123-I1): لا يلزم أن يكون c من مخزون الثوابت المميزة؛ نجدد ثوابت البرهان ثم نبدله بثابت من E.
خذ برهانًا منتهيًا~$\Pi$ لـ~$!A(c)$ من~$\Gamma$، وسمّ فروضه المستعملة
$\Gamma_0$. جدّد ثوابته المميَّزة بـ
  \olref[qua]{lem:qr-freshening} مع مجموعة المنع~$\{c\}$، ثم اختر
$e\in\mathcal E\setminus\{c\}$ لا يقع في شيء من صيغ البرهان المجدَّد ولا يساوي ثابتًا
مميَّزًا مكتوبًا على عقدة~\QR{} فيه. بدّل~$c$ و~$e$ تقابليًا في البرهان
كله. ولا يتغير~$\Gamma_0$، إذ كلا الثابتين غير وارد فيه؛ كما تحفظ إعادة
التسمية البديهيات وإثبات المقدَّم وصورتي~\QR{} وشروطهما. فينتج
$\Gamma_0\Proves !A(e)$.

  ونضم إلى ذلك الصيغة~$!A(e)\lif(\ltrue\lif !A(e))$ من بديهية اللزوم
  الأولى، فنستخرج بإثبات المقدَّم
  $\Gamma_0\Proves\ltrue\lif !A(e)$. ولما كان~$e\in\mathcal E$ غير وارد
في~$\Gamma_0$ ولا~$\ltrue$ ولا~$!A(x)$، جاز إجراء~\QR{}، فحصلنا على
$\Gamma_0\Proves\ltrue\lif\lforall[x][!A(x)]$. ونضم إليه
$\Gamma_0\Proves\ltrue$، فنستنتج بإثبات المقدَّم
$\Gamma_0\Proves\lforall[x][!A(x)]$؛ ثم ترفع الرتابة النتيجة إلى~$\Gamma$.
\end{proof}

\begin{prop}
\ollabel{prop:provability-quantifiers}
لكل حد مغلق~$t$ تصح العبارتان الآتيتان:
\begin{tagenumerate}{prvEx,prvAll}
\tagitem{prvEx}{$!A(t) \Proves \lexists[x][!A(x)]$.}{}

\tagitem{prvAll}{$\lforall[x][!A(x)] \Proves !A(t)$.}{}
\end{tagenumerate}
\end{prop}

\begin{proof}
\begin{tagenumerate}{prvEx,prvAll}
\tagitem{prvEx}{نأخذ \olref[qua]{ax:q2} ثم نطبق مبرهنة الاستنباط.}{}
\tagitem{prvAll}{ونأخذ \olref[qua]{ax:q1} ثم نطبق مبرهنة الاستنباط.}{}
\end{tagenumerate}
\end{proof}

\end{document}
